\chapter{Sprint 2: Rapport Request and Form Templates}

\section{Introduction}

Sprint 2 implemented the first business-critical capabilities of RM by introducing intervention request management and dynamic form template governance. This sprint transformed the platform from an authenticated infrastructure into an operational system that can create, assign, and process intervention activities with structured data capture.

The sprint objective was to digitize the intervention initiation process and standardize data collection through configurable templates.

\section{Preview of Sprint}

\subsection{Sprint Objectives}

\begin{enumerate}
\item Implement intervention request lifecycle creation and assignment.
\item Implement rapport template management with JSON mapping definitions.
\item Build dynamic form generation rules from template configuration.
\item Enforce validation constraints for reliable and consistent submissions.
\item Expose API contracts for web and mobile interaction with intervention data.
\end{enumerate}

\subsection{Business Value Statement}

This sprint addresses the PRD objective of replacing fragmented request processes with a structured digital workflow. It provides standardized form definitions, reducing ambiguity and improving data quality before validation and reporting phases.

\section{Sprint Backlog}

\begin{table}[H]
\centering
\caption{Sprint 2 Backlog}
\label{tab:sprint2_backlog}
\begin{tabularx}{\textwidth}{|l|X|l|l|}
\hline
\textbf{ID} & \textbf{User Story} & \textbf{Points} & \textbf{Status} \\
\hline
S2-1 & Create and assign rapport requests to technicians & 13 & Completed \\
\hline
S2-2 & Configure and maintain rapport templates & 8 & Completed \\
\hline
S2-3 & Generate dynamic forms from template mappings & 8 & Completed \\
\hline
S2-4 & Validate input data and enforce required fields & 8 & Completed \\
\hline
S2-5 & Support request filtering and lookup in web interface & 8 & Completed \\
\hline
\textbf{Total} & & \textbf{45} & \\
\hline
\end{tabularx}
\end{table}

\section{Functional Specifications}

\subsection{Detailed Use Cases}

\textbf{UC-S2-01: Create Rapport Request}

\textbf{Primary Actor:} Project Coordinator

\textbf{Preconditions:}
\begin{itemize}
\item Coordinator is authenticated and authorized.
\item A valid template and client context exist.
\end{itemize}

\textbf{Main Flow:}
\begin{enumerate}
\item Coordinator selects project and client context.
\item Coordinator chooses rapport template.
\item Coordinator enters planning attributes and assignment details.
\item System validates mandatory metadata.
\item System creates request with initial status \texttt{PLANNED}.
\item Technician receives assignment visibility.
\end{enumerate}

\textbf{Postconditions:}
\begin{itemize}
\item Intervention request exists with traceable metadata and ownership.
\item Request can be consumed by mobile execution workflow.
\end{itemize}

\textbf{UC-S2-02: Configure Rapport Template}

\textbf{Primary Actor:} Project Coordinator

\textbf{Main Flow:}
\begin{enumerate}
\item Coordinator creates template descriptor (name, project, delay policy).
\item Coordinator uploads or edits mapping JSON definition.
\item System validates schema consistency.
\item Template is persisted and marked active.
\item Generated forms become available for requests linked to the template.
\end{enumerate}

\textbf{UC-S2-03: Render Dynamic Form on Mobile}

\textbf{Primary Actor:} Field Technician

\textbf{Main Flow:}
\begin{enumerate}
\item Technician opens assigned intervention.
\item Mobile client fetches linked template mapping.
\item Form engine renders controls by field type.
\item Validation rules are attached to each field.
\item Technician fills form progressively and saves state.
\end{enumerate}

\subsection{Textual Specifications}

\begin{table}[H]
\centering
\caption{Sprint 2 Functional Specification Summary}
\label{tab:sprint2_specs}
\begin{tabularx}{\textwidth}{|l|X|X|}
\hline
\textbf{Spec ID} & \textbf{Specification} & \textbf{Acceptance Rule} \\
\hline
S2-FS-01 & Request creation with template linkage & Request stores template reference and initial status \\
\hline
S2-FS-02 & Template CRUD operations for coordinator role & Unauthorized roles cannot modify template resources \\
\hline
S2-FS-03 & Dynamic form rendering from mapping JSON & Field set is generated without manual UI coding per template \\
\hline
S2-FS-04 & Required and typed field validation & Submission blocked when constraints are violated \\
\hline
S2-FS-05 & Paginated filtering for intervention search & Query by status/client/reference returns consistent results \\
\hline
\end{tabularx}
\end{table}

\section{Conception}

\subsection{Sequence Diagram (Placeholder)}

\begin{figure}[H]
\centering
\fbox{\parbox{0.9\textwidth}{
\centering
% TODO: Insert Sprint 2 sequence diagram here
Sequence Diagram Placeholder: Coordinator request creation, template retrieval, request persistence, technician retrieval, dynamic rendering.
}}
\caption{Sprint 2 Request and Template Sequence Placeholder}
\label{fig:sprint2_sequence_placeholder}
\end{figure}

\subsection{Class Diagram (Placeholder)}

\begin{figure}[H]
\centering
\fbox{\parbox{0.9\textwidth}{
\centering
% TODO: Insert Sprint 2 class diagram here
Class Diagram Placeholder: RapportRequest, RapportTemplate, Project, Client, ReportLineData, and validation rule classes.
}}
\caption{Sprint 2 Domain Class Diagram Placeholder}
\label{fig:sprint2_class_placeholder}
\end{figure}

\section{Realization}

\subsection{Backend Implementation Details (Spring Boot)}

Backend realization in this sprint focused on the intervention and template domains:

\begin{itemize}
\item request creation and update endpoints,
\item status initialization and transition safeguards,
\item template persistence with mapping JSON,
\item validation rules at DTO and service layers,
\item filtering and pagination specifications for operational dashboards.
\end{itemize}

Primary endpoint groups delivered:

\begin{itemize}
\item \api{rapport-requests}
\item \api{rapport-requests-custom}
\item \api{rapport-templates}
\end{itemize}

\subsection{Frontend Implementation (Angular)}

The Angular web application delivered:

\begin{itemize}
\item intervention list with filters by status, client, and reference,
\item request creation forms for coordinators,
\item template management screens (create/edit/delete),
\item validation feedback for invalid input states.
\end{itemize}

\subsection{Mobile Implementation (Flutter)}

The Flutter mobile realization delivered:

\begin{itemize}
\item assigned intervention retrieval with pagination,
\item dynamic form rendering from template mapping,
\item local form state handling before final submission,
\item client-side validation for mandatory fields and field format rules.
\end{itemize}

\subsection{Database Design (PostgreSQL)}

Sprint 2 extended the schema with business entities required by intervention workflows:

\begin{itemize}
\item \texttt{rapport\_request} table for intervention records,
\item \texttt{rapport\_template} table for form definitions,
\item JSONB columns for flexible template mapping and dynamic line data,
\item foreign key relationships to project, client, and assignee contexts.
\end{itemize}

\begin{figure}[H]
\centering
\fbox{\parbox{0.88\textwidth}{
\centering
% TODO: Insert Sprint 2 database relational diagram here
Database Placeholder: Rapport request and template schema with project/client/user relations.
}}
\caption{Sprint 2 Database Schema Placeholder}
\label{fig:sprint2_db_placeholder}
\end{figure}

\subsection{API and Endpoint Summary}

\begin{table}[H]
\centering
\caption{Sprint 2 API Summary}
\label{tab:sprint2_api_summary}
\begin{tabularx}{\textwidth}{|l|l|X|}
\hline
\textbf{Method} & \textbf{Endpoint} & \textbf{Purpose} \\
\hline
POST & \api{rapport-requests} & Create intervention request from template metadata \\
\hline
GET & \api{rapport-requests} & Retrieve paginated and filtered intervention list \\
\hline
GET & \api{rapport-requests/{id}} & Retrieve intervention details for review and execution \\
\hline
POST & \api{rapport-templates} & Create a new dynamic form template \\
\hline
PUT & \api{rapport-templates} & Update template configuration and mapping JSON \\
\hline
DELETE & \api{rapport-templates/{id}} & Archive or remove template according to policy \\
\hline
\end{tabularx}
\end{table}

\section{Conclusion}

Sprint 2 transformed RM into an operational intervention platform by delivering request lifecycle management and configurable dynamic forms. It established the data quality baseline for downstream approval, evidence management, and reporting functionalities. Sprint 3 builds on this core by adding file management and geolocation tracking to strengthen field execution traceability.
